\documentclass[a4paper, 12pt]{article}

\usepackage{utils}

\renewcommand*{\today}{27 January 2026}

\newtcbtheorem[number within=section]{recherche}{Recherche}%
{
    enhanced,
    colframe=black,
    coltitle=black, fonttitle=\bfseries,
    top=5mm,
    attach boxed title to top left={yshift=-3mm, xshift=6mm},
    boxed title style={colback=white, arc=2mm, shadow={1mm}{-1mm}{0mm}{black}}
}{recherche}

\begin{document}

\hotbox{Questions ouvertes}{CM 2}{\today}

\section{27 Janvier 2026}

\subsection*{Travaux réalisés}

\begin{itemize}[label=\textbullet]
    \item Amélioration du programme (prise en compte des multipliant supérieurs à 9) (\textbf{Annexe A}).
    \item Modélisation, en posant $y := \frac{x - d}{10}$ (le nombre tournant sans son chiffre des unités), on obtient l'équation centrale suivante : $y = \frac{d(m - 10^{l - 1})}{1 - 10m}$.
\end{itemize}

\subsection*{Nouvelles questions}

\begin{itemize}[label=\textbullet]
    \item \textbf{[Q2]} Comment trouver la longueur $l$ d'un nombre tournant à partir de son multipliant $m$ et de son chiffre des unités $d$ ?
\end{itemize}

\end{document}